% Adaptación de FPC/modules/fcp_gg.tex. No contiene series nuevas.
\begin{frame}{¿Qué entendemos por ingresos y gastos?}
\begin{itemize}
  \item No todo dinero que entra es un ingreso fiscal.
  \item No todo dinero que sale es un gasto fiscal.
  \item El presupuesto autoriza operaciones que, en las cuentas fiscales, pueden ser \textbf{financiamiento}.
\end{itemize}
\vspace{1em}
La pregunta es qué operación económica hay detrás del movimiento de caja.
\fuenteFMI
\end{frame}

\begin{frame}{Recordatorio contable: ¿qué es el patrimonio?}
\[\text{Patrimonio neto}=\text{Activos}-\text{Pasivos}\]
\begin{itemize}
  \item \textbf{Activos}: efectivo, cuentas por cobrar, acciones, infraestructura.
  \item \textbf{Pasivos}: préstamos, bonos, cuentas por pagar.
  \item Un intercambio entre activos o el pago del principal de una deuda no crea, por sí solo, riqueza.
\end{itemize}
\fuenteFMI
\end{frame}

\begin{frame}{El ingreso fiscal no incluye el endeudamiento}
\begin{itemize}
  \item En el MEFP, los \textbf{ingresos} son transacciones que aumentan el patrimonio neto: impuestos, contribuciones, transferencias recibidas y rentas de la propiedad, entre otras.
  \item El crédito recibido aumenta simultáneamente la caja y la deuda.
  \item Por eso el crédito \textbf{financia} el déficit; no lo reduce como lo haría un impuesto.
\end{itemize}
\fuenteFMI
\end{frame}

\begin{frame}{¿Qué se incluye en el gasto fiscal?}
\begin{itemize}
  \item En el MEFP, el gasto que reduce patrimonio incluye salarios, intereses, subsidios y transferencias, entre otros conceptos.
  \item Para calcular el balance fiscal también se considera la \textbf{inversión neta en activos no financieros}, como infraestructura y equipos.
  \item El pago del principal de la deuda es una operación de \textbf{financiamiento}, no gasto que aumente el déficit.
\end{itemize}
\vspace{0.5em}
\small No confundir el resultado operativo con el balance fiscal: construir un activo también requiere recursos.
\fuenteFMI
\end{frame}

\begin{frame}{¿Pedir prestados \$100 millones es un ingreso?}
\begin{align*}
  \Delta\text{Efectivo}&=+100\\
  \Delta\text{Deuda}&=+100\\
  \Delta\text{Patrimonio neto}&=100-100=0
\end{align*}
\begin{itemize}
  \item Hay más liquidez, pero también una obligación de devolución.
  \item \textbf{No hay ingreso fiscal}: hay financiamiento.
\end{itemize}
\fuenteFMI
\end{frame}

\begin{frame}{¿Por qué las amortizaciones no son gasto fiscal?}
Pagar \$20 millones del principal de una deuda:
\begin{align*}
  \Delta\text{Efectivo}&=-20\\
  \Delta\text{Deuda}&=-20\\
  \Delta\text{Patrimonio neto}&=-20-(-20)=0
\end{align*}
\textbf{La caja disminuye, pero la deuda también.}
\fuenteFMI
\end{frame}

\begin{frame}{Los intereses sí son gasto}
\begin{itemize}
  \item Remuneran el uso de recursos prestados.
  \item Su causación reduce el patrimonio sin amortizar el principal.
  \item El servicio presupuestal de la deuda incluye tanto \textbf{intereses} como \textbf{amortizaciones}; su tratamiento fiscal es distinto.
\end{itemize}
\fuenteFMI
\fuenteEOP{36}
\end{frame}

\begin{frame}{¿Qué parte de una cuota hipotecaria es gasto?}
Pago mensual: \textbf{\$6 millones}.
\vspace{0.8em}
\begin{center}
\begin{tabular}{lrr}
\toprule
Concepto & Salida de caja & Gasto fiscal\\
\midrule
Intereses & 4 & 4\\
Amortización & 2 & 0\\
\midrule
Total & 6 & 4\\
\bottomrule
\end{tabular}
\end{center}
\vspace{0.6em}
Los \$2 millones restantes reducen la deuda, no el balance fiscal.
\fuenteFMI
\end{frame}

\begin{frame}{¿Vender un apartamento es un ingreso?}
Se vende por \$500 millones un apartamento registrado por ese valor.
\begin{itemize}
  \item Aumenta el efectivo en 500 y disminuye otro activo en 500.
  \item No aumenta el patrimonio por la transacción.
  \item En el MEFP no es ingreso: es \textbf{venta de un activo no financiero}.
\end{itemize}
\vspace{0.7em}
\textbf{Pero sí afecta el balance de préstamo neto/endeudamiento neto}: reduce la inversión neta en activos no financieros.
\fuenteFMI
\end{frame}

\begin{frame}{¿Y vender acciones de una empresa?}
\begin{itemize}
  \item Se intercambia un \textbf{activo financiero} por otro: acciones por efectivo.
  \item En el MEFP, la transacción se registra como financiamiento, no como ingreso.
  \item A diferencia del apartamento, no modifica la inversión neta en activos \textbf{no financieros}.
\end{itemize}
\vspace{0.7em}
La clasificación del activo cambia el tratamiento fiscal de su venta.
\fuenteFMI
\end{frame}

\begin{frame}{La venta de ISA permite plantear la pregunta}
En 2021, la Nación vendió a Ecopetrol su participación en ISA por aproximadamente \textbf{\$14,2 billones}.
\begin{itemize}
  \item La operación aportó caja y cambió la tenencia de activos.
  \item No produjo una nueva fuente permanente de recaudo.
  \item Para evaluar su efecto fiscal hay que preguntar: ¿con qué metodología se registra y qué entidades se consolidan?
\end{itemize}
\note{[Sources] Ejemplo adaptado de FPC, modules/fcp_gg.tex. MHCP, Estabilización de la deuda y protección de los recursos públicos, 29 de abril de 2022: https://www.minhacienda.gov.co/w/estabilizacion-de-la-deuda-y-proteccion-de-los-recursos-publicos-entre-los-logros-del-sector-hacienda-a-100-dias-de-cierre-de-gobierno. Se separa el criterio del MEFP del registro histórico.}
\end{frame}

\begin{frame}{Balance fiscal y déficit}
\[B=T-G \qquad \text{Déficit}=-B\]
\begin{itemize}
  \item $T$: ingresos fiscales; $G$: gasto fiscal total, incluidos intereses.
  \item Si $B<0$, hay déficit. Si $B>0$, hay superávit.
  \item En esta identidad simplificada no hay préstamo neto ni otros ajustes.
\end{itemize}
\vspace{0.6em}
Con ingresos de 80 y gasto de 100, el balance es $-20$ y el déficit es 20.
\fuenteFMI
\fuenteCaja
\end{frame}

\begin{frame}{¿Cómo se pasa a porcentaje del PIB?}
\[b=100\frac{B}{PIB},\qquad t=100\frac{T}{PIB},\qquad g=100\frac{G}{PIB}\]
\[b=t-g\]
\begin{itemize}
  \item Si el PIB es 1.000 y el balance es $-20$, el balance es $-2\%$ del PIB.
  \item Un punto del PIB equivale a 10 en este ejemplo.
\end{itemize}
\fuenteFMI
\end{frame}

\begin{frame}{Gasto primario y balance primario}
\begin{align*}GP &=G-I\\ P &=T-GP=B+I\end{align*}
\begin{itemize}
  \item El gasto primario excluye los intereses $I$.
  \item El balance primario muestra el resultado antes de intereses.
  \item Con $T=80$, $G=100$ e $I=15$: $GP=85$ y $P=-5$.
\end{itemize}
\fuenteFMI
\end{frame}

\begin{frame}{El déficit no es toda la necesidad de caja}
Ejemplo sin variación de activos financieros ni otros ajustes:
\begin{align*}
  \text{Déficit}&=20\\
  \text{Amortizaciones}&=30\\
  \text{Financiamiento bruto requerido}&=20+30=50
\end{align*}
\textbf{Emitir deuda por 50 no significa tener un déficit de 50.}
\fuenteFMI
\end{frame}

\begin{frame}{Presupuesto y cuentas fiscales responden preguntas distintas}
\begin{itemize}
  \item \textbf{Presupuesto}: ¿qué recursos se estiman y qué operaciones se autoriza realizar?
  \item \textbf{Ejecución presupuestal}: ¿qué se comprometió, obligó, pagó y recaudó?
  \item \textbf{Cuentas fiscales}: ¿qué ingresos, gastos y financiamiento resultan bajo una metodología y cobertura definidas?
\end{itemize}
\vspace{0.7em}
La contabilidad fiscal no se obtiene copiando el total del PGN.
\fuenteEOP{7, 11, 31 y 36}
\fuenteCaja
\end{frame}
